\documentclass[a4paper]{article}

%%%%%%%%%%%%%%%%%%%%%%%%%%%%%%%%%%%%%%%%%%%%%%%%%%%%%%%%%%%%%%%%%%%%%%%%%%%%%%%

\usepackage[utf8]{inputenc}
\usepackage[T1]{fontenc}
\usepackage[english]{babel}
\usepackage{caption}
\usepackage{subcaption}
\usepackage{float}
\usepackage{hyperref}
\usepackage{fancyhdr}
\usepackage[a4paper, left=2.85cm, right=2.85cm]{geometry}
\usepackage{graphicx}
\usepackage[nottoc]{tocbibind}
\usepackage[bottom]{footmisc}
\graphicspath{{images/}}

\pagestyle{fancy}
\fancyhf{}
\rhead{\small Paris-Saclay University -- 23/12/2021}
\lhead{\small Gaëtan Serré -- Master AI}
\rfoot{\thepage}

\begin{document}

\newgeometry{
  bottom=1.5cm
}
\setlength{\headheight}{13.59999pt}
\addtolength{\topmargin}{-1.59999pt}

\begin{center}
  \vspace*{5cm}

  \textbf{\large Learning to Run a Power Network Energies of the future and carbon neutrality}

  \vspace*{1cm}

  Gaëtan Serré \\
  {\ttfamily gaetan.serre@universite-paris-saclay.fr}

  \vspace*{10cm}

  \vfill
  \noindent\rule{\textwidth}{0.4pt}
  \begin{figure}[H]
      \centering
      \begin{subfigure}{.35\textwidth}
          \includegraphics[width=\textwidth]{ups-logo.png}
      \end{subfigure}
  \end{figure}
\end{center}
\pagenumbering{gobble}

\pagebreak
\restoregeometry

\pagenumbering{arabic}

\section{Algorithm \& Results}
The goal of this exercise was to provide a reinforcement learning solution to the challenge
\textit{L2RPN NEURIPS 2020}\cite{L2RPN}. I chose to implement a Deep Q-Learning algorithm.
For reasons of time and computational power, I was able to train my model during 2,000 episodes
considering only 50 actions (including the \textit{do-nothing} action). I used the default reward.
We can see on \ref{table:results} that the more the number of episodes increases, the longer the model survives
\begin{table}[H]
  \centering
  \begin{tabular}{|l|c|r|}
    \hline
    Episode & Average steps survived & Average total reward \\
    \hline
    0-100 & 16 & 6,000 \\
    \hline
    100-300 & 32 & 20,000 \\
    \hline
    400-500 & 100 & 100,000\\
    \hline
    600-700 & 500 & 500,000\\
    \hline
    800-1000 & 900 & 900,000\\
    \hline
    1000-1500 & 1200 & 1,000,000\\
    \hline
    1500-2000 & 2000 & 1,800,000\\
    \hline
  \end{tabular}
  \caption{Summary of the results of the training phase}
  \label{table:results}
\end{table}

My model reaches a score of $3.423101$ on the competition, which is slightly better than the \textit{do-nothing}
strategy. I believe that, with more episodes and some optimizations on the neural network
used or on the selection features of the power grid, my algorithm might have a better score.

\section{Difficulties}
During these few days, I encountered some difficulties.
\begin{itemize}
  \item The first one was first of all
        understand the purpose of the challenge and what was the stake of the problem 
        because I was not familiar with power grid.
        The starting kit and the paper
        \textit{Learning to run a Power Network Challenge: a Retrospective Analysis}
        \cite{marot2021learning} helped me a lot.
  \item Another difficulty was to use the \textit{grid2op} library.
        I had to do a lot of testing to find out what was available to me
        to make a reinforcement learning model. The starting kit also helped me a lot.
  \item The choice of the model was quite natural, I had already implemented a Q-Learning
        algorithm and I thought that Deep Q-Learning was very suitable for the problem.
        However, I had never implemented it before so I had to follow some tutorials.
  \item Features selection was also a problem. Beyond a version error of \textit{grid2op}
        that gave me a hard time,
        looking at the public submissions, I could see that the best teams were selecting features
        of the power grid. As I don't have the expert knowledge to do the same,
        I decided to keep them all but it is surely not the most optimized solution.
  \item The space action having more than 70,000 possible actions, I had to select only a subset of them
        in order to save computing time. I was able to use the file \textit{top1000\_actions.npz} provided
        by the winning team \cite{Zhou2021ActionSB}. This file contains the 1,000 best actions according to them.
        I chose to keep only the first 50 while being sure that the \textit{do-nothing} action is included.
        This way, I knew that I could never have the best solution but I could do as good
        as the \textit{do-nothing} agent.
  \item The last difficulty was to participate in a more "professional" competition
        than the ones I had participated in.
        More specifically, formatting my submissions took a lot of time.

\end{itemize}

\section{What I have learned}
By surpassing the difficulties mentioned above, I learned to implement a
Deep Q-Learning algorithm and to participate in a professional competition. Thanks to
\textit{grid2op}, I was able to learn the vocabulary and interfaces
of reinforcement learning. My knowledge in this domain came from
personal projects in which I implemented everything \textit{from scratch},
so I didn't have this general culture. Finally, I had an overview of the problem of managing
power grids and what is at stake for our society.

\section{Conclusion}
Before I even started, I was already very motivated by this internship.
I learned about reinforcement learning through its application
in game theory and I have been very interested in it since then.
Thanks to this challenge, I discovered a more "realistic" application than Go
or chess of reinforcement learning and I really liked it.
I would like to go deeper into the problem by improving my algorithm or
or by trying others.

\bibliographystyle{abbrv}
\bibliography{refs}


\end{document}